% Options for packages loaded elsewhere
% Options for packages loaded elsewhere
\PassOptionsToPackage{unicode}{hyperref}
\PassOptionsToPackage{hyphens}{url}
\PassOptionsToPackage{dvipsnames,svgnames,x11names}{xcolor}
%
\documentclass[
  letterpaper,
  DIV=11,
  numbers=noendperiod]{scrartcl}
\usepackage{xcolor}
\usepackage{amsmath,amssymb}
\setcounter{secnumdepth}{-\maxdimen} % remove section numbering
\usepackage{iftex}
\ifPDFTeX
  \usepackage[T1]{fontenc}
  \usepackage[utf8]{inputenc}
  \usepackage{textcomp} % provide euro and other symbols
\else % if luatex or xetex
  \usepackage{unicode-math} % this also loads fontspec
  \defaultfontfeatures{Scale=MatchLowercase}
  \defaultfontfeatures[\rmfamily]{Ligatures=TeX,Scale=1}
\fi
\usepackage{lmodern}
\ifPDFTeX\else
  % xetex/luatex font selection
\fi
% Use upquote if available, for straight quotes in verbatim environments
\IfFileExists{upquote.sty}{\usepackage{upquote}}{}
\IfFileExists{microtype.sty}{% use microtype if available
  \usepackage[]{microtype}
  \UseMicrotypeSet[protrusion]{basicmath} % disable protrusion for tt fonts
}{}
\makeatletter
\@ifundefined{KOMAClassName}{% if non-KOMA class
  \IfFileExists{parskip.sty}{%
    \usepackage{parskip}
  }{% else
    \setlength{\parindent}{0pt}
    \setlength{\parskip}{6pt plus 2pt minus 1pt}}
}{% if KOMA class
  \KOMAoptions{parskip=half}}
\makeatother
% Make \paragraph and \subparagraph free-standing
\makeatletter
\ifx\paragraph\undefined\else
  \let\oldparagraph\paragraph
  \renewcommand{\paragraph}{
    \@ifstar
      \xxxParagraphStar
      \xxxParagraphNoStar
  }
  \newcommand{\xxxParagraphStar}[1]{\oldparagraph*{#1}\mbox{}}
  \newcommand{\xxxParagraphNoStar}[1]{\oldparagraph{#1}\mbox{}}
\fi
\ifx\subparagraph\undefined\else
  \let\oldsubparagraph\subparagraph
  \renewcommand{\subparagraph}{
    \@ifstar
      \xxxSubParagraphStar
      \xxxSubParagraphNoStar
  }
  \newcommand{\xxxSubParagraphStar}[1]{\oldsubparagraph*{#1}\mbox{}}
  \newcommand{\xxxSubParagraphNoStar}[1]{\oldsubparagraph{#1}\mbox{}}
\fi
\makeatother

\usepackage{color}
\usepackage{fancyvrb}
\newcommand{\VerbBar}{|}
\newcommand{\VERB}{\Verb[commandchars=\\\{\}]}
\DefineVerbatimEnvironment{Highlighting}{Verbatim}{commandchars=\\\{\}}
% Add ',fontsize=\small' for more characters per line
\usepackage{framed}
\definecolor{shadecolor}{RGB}{241,243,245}
\newenvironment{Shaded}{\begin{snugshade}}{\end{snugshade}}
\newcommand{\AlertTok}[1]{\textcolor[rgb]{0.68,0.00,0.00}{#1}}
\newcommand{\AnnotationTok}[1]{\textcolor[rgb]{0.37,0.37,0.37}{#1}}
\newcommand{\AttributeTok}[1]{\textcolor[rgb]{0.40,0.45,0.13}{#1}}
\newcommand{\BaseNTok}[1]{\textcolor[rgb]{0.68,0.00,0.00}{#1}}
\newcommand{\BuiltInTok}[1]{\textcolor[rgb]{0.00,0.23,0.31}{#1}}
\newcommand{\CharTok}[1]{\textcolor[rgb]{0.13,0.47,0.30}{#1}}
\newcommand{\CommentTok}[1]{\textcolor[rgb]{0.37,0.37,0.37}{#1}}
\newcommand{\CommentVarTok}[1]{\textcolor[rgb]{0.37,0.37,0.37}{\textit{#1}}}
\newcommand{\ConstantTok}[1]{\textcolor[rgb]{0.56,0.35,0.01}{#1}}
\newcommand{\ControlFlowTok}[1]{\textcolor[rgb]{0.00,0.23,0.31}{\textbf{#1}}}
\newcommand{\DataTypeTok}[1]{\textcolor[rgb]{0.68,0.00,0.00}{#1}}
\newcommand{\DecValTok}[1]{\textcolor[rgb]{0.68,0.00,0.00}{#1}}
\newcommand{\DocumentationTok}[1]{\textcolor[rgb]{0.37,0.37,0.37}{\textit{#1}}}
\newcommand{\ErrorTok}[1]{\textcolor[rgb]{0.68,0.00,0.00}{#1}}
\newcommand{\ExtensionTok}[1]{\textcolor[rgb]{0.00,0.23,0.31}{#1}}
\newcommand{\FloatTok}[1]{\textcolor[rgb]{0.68,0.00,0.00}{#1}}
\newcommand{\FunctionTok}[1]{\textcolor[rgb]{0.28,0.35,0.67}{#1}}
\newcommand{\ImportTok}[1]{\textcolor[rgb]{0.00,0.46,0.62}{#1}}
\newcommand{\InformationTok}[1]{\textcolor[rgb]{0.37,0.37,0.37}{#1}}
\newcommand{\KeywordTok}[1]{\textcolor[rgb]{0.00,0.23,0.31}{\textbf{#1}}}
\newcommand{\NormalTok}[1]{\textcolor[rgb]{0.00,0.23,0.31}{#1}}
\newcommand{\OperatorTok}[1]{\textcolor[rgb]{0.37,0.37,0.37}{#1}}
\newcommand{\OtherTok}[1]{\textcolor[rgb]{0.00,0.23,0.31}{#1}}
\newcommand{\PreprocessorTok}[1]{\textcolor[rgb]{0.68,0.00,0.00}{#1}}
\newcommand{\RegionMarkerTok}[1]{\textcolor[rgb]{0.00,0.23,0.31}{#1}}
\newcommand{\SpecialCharTok}[1]{\textcolor[rgb]{0.37,0.37,0.37}{#1}}
\newcommand{\SpecialStringTok}[1]{\textcolor[rgb]{0.13,0.47,0.30}{#1}}
\newcommand{\StringTok}[1]{\textcolor[rgb]{0.13,0.47,0.30}{#1}}
\newcommand{\VariableTok}[1]{\textcolor[rgb]{0.07,0.07,0.07}{#1}}
\newcommand{\VerbatimStringTok}[1]{\textcolor[rgb]{0.13,0.47,0.30}{#1}}
\newcommand{\WarningTok}[1]{\textcolor[rgb]{0.37,0.37,0.37}{\textit{#1}}}

\usepackage{longtable,booktabs,array}
\usepackage{calc} % for calculating minipage widths
% Correct order of tables after \paragraph or \subparagraph
\usepackage{etoolbox}
\makeatletter
\patchcmd\longtable{\par}{\if@noskipsec\mbox{}\fi\par}{}{}
\makeatother
% Allow footnotes in longtable head/foot
\IfFileExists{footnotehyper.sty}{\usepackage{footnotehyper}}{\usepackage{footnote}}
\makesavenoteenv{longtable}
\usepackage{graphicx}
\makeatletter
\newsavebox\pandoc@box
\newcommand*\pandocbounded[1]{% scales image to fit in text height/width
  \sbox\pandoc@box{#1}%
  \Gscale@div\@tempa{\textheight}{\dimexpr\ht\pandoc@box+\dp\pandoc@box\relax}%
  \Gscale@div\@tempb{\linewidth}{\wd\pandoc@box}%
  \ifdim\@tempb\p@<\@tempa\p@\let\@tempa\@tempb\fi% select the smaller of both
  \ifdim\@tempa\p@<\p@\scalebox{\@tempa}{\usebox\pandoc@box}%
  \else\usebox{\pandoc@box}%
  \fi%
}
% Set default figure placement to htbp
\def\fps@figure{htbp}
\makeatother





\setlength{\emergencystretch}{3em} % prevent overfull lines

\providecommand{\tightlist}{%
  \setlength{\itemsep}{0pt}\setlength{\parskip}{0pt}}



 


\KOMAoption{captions}{tableheading}
\makeatletter
\@ifpackageloaded{caption}{}{\usepackage{caption}}
\AtBeginDocument{%
\ifdefined\contentsname
  \renewcommand*\contentsname{Table of contents}
\else
  \newcommand\contentsname{Table of contents}
\fi
\ifdefined\listfigurename
  \renewcommand*\listfigurename{List of Figures}
\else
  \newcommand\listfigurename{List of Figures}
\fi
\ifdefined\listtablename
  \renewcommand*\listtablename{List of Tables}
\else
  \newcommand\listtablename{List of Tables}
\fi
\ifdefined\figurename
  \renewcommand*\figurename{Figure}
\else
  \newcommand\figurename{Figure}
\fi
\ifdefined\tablename
  \renewcommand*\tablename{Table}
\else
  \newcommand\tablename{Table}
\fi
}
\@ifpackageloaded{float}{}{\usepackage{float}}
\floatstyle{ruled}
\@ifundefined{c@chapter}{\newfloat{codelisting}{h}{lop}}{\newfloat{codelisting}{h}{lop}[chapter]}
\floatname{codelisting}{Listing}
\newcommand*\listoflistings{\listof{codelisting}{List of Listings}}
\makeatother
\makeatletter
\makeatother
\makeatletter
\@ifpackageloaded{caption}{}{\usepackage{caption}}
\@ifpackageloaded{subcaption}{}{\usepackage{subcaption}}
\makeatother
\usepackage{bookmark}
\IfFileExists{xurl.sty}{\usepackage{xurl}}{} % add URL line breaks if available
\urlstyle{same}
\hypersetup{
  pdftitle={Matplotlib},
  pdfauthor={Alan Lujan},
  colorlinks=true,
  linkcolor={blue},
  filecolor={Maroon},
  citecolor={Blue},
  urlcolor={Blue},
  pdfcreator={LaTeX via pandoc}}


\title{Matplotlib}
\usepackage{etoolbox}
\makeatletter
\providecommand{\subtitle}[1]{% add subtitle to \maketitle
  \apptocmd{\@title}{\par {\large #1 \par}}{}{}
}
\makeatother
\subtitle{Data Visualization for Economic Analysis}
\author{Alan Lujan}
\date{}
\begin{document}
\maketitle


\subsection{What is Matplotlib?}\label{what-is-matplotlib}

\textbf{Matplotlib} is Python's primary plotting library for creating
publication-quality figures.

\textbf{Why visualization matters in economics:} - Communicate results
effectively - Explore data patterns - Present model predictions vs.~data
- Diagnostic plots for model checking - Publication-ready charts for
papers

\textbf{Interface:} MATLAB-like API + object-oriented interface

\subsection{Basic Plotting Workflow}\label{basic-plotting-workflow}

\textbf{Simple plot:}

\begin{Shaded}
\begin{Highlighting}[]
\ImportTok{import}\NormalTok{ matplotlib.pyplot }\ImportTok{as}\NormalTok{ plt}
\ImportTok{import}\NormalTok{ numpy }\ImportTok{as}\NormalTok{ np}

\NormalTok{x }\OperatorTok{=}\NormalTok{ np.linspace(}\DecValTok{0}\NormalTok{, }\DecValTok{10}\NormalTok{, }\DecValTok{100}\NormalTok{)}
\NormalTok{y }\OperatorTok{=}\NormalTok{ np.sin(x)}

\NormalTok{plt.plot(x, y)}
\NormalTok{plt.xlabel(}\StringTok{\textquotesingle{}x\textquotesingle{}}\NormalTok{)}
\NormalTok{plt.ylabel(}\StringTok{\textquotesingle{}sin(x)\textquotesingle{}}\NormalTok{)}
\NormalTok{plt.title(}\StringTok{\textquotesingle{}Sine Function\textquotesingle{}}\NormalTok{)}
\NormalTok{plt.show()}
\end{Highlighting}
\end{Shaded}

\textbf{Result:} A line plot with labeled axes and title.

\subsection{Figure and Axes Objects}\label{figure-and-axes-objects}

\textbf{Object-oriented approach} (recommended):

\begin{Shaded}
\begin{Highlighting}[]
\NormalTok{fig, ax }\OperatorTok{=}\NormalTok{ plt.subplots()  }\CommentTok{\# Create figure and axes}

\NormalTok{ax.plot(x, y)}
\NormalTok{ax.set\_xlabel(}\StringTok{\textquotesingle{}x\textquotesingle{}}\NormalTok{)}
\NormalTok{ax.set\_ylabel(}\StringTok{\textquotesingle{}sin(x)\textquotesingle{}}\NormalTok{)}
\NormalTok{ax.set\_title(}\StringTok{\textquotesingle{}Sine Function\textquotesingle{}}\NormalTok{)}

\NormalTok{plt.show()}
\end{Highlighting}
\end{Shaded}

\textbf{Why OO approach?} - More control over multiple subplots -
Cleaner code for complex figures - Better for functions that return
plots

\subsection{Multiple Lines and
Legends}\label{multiple-lines-and-legends}

\begin{Shaded}
\begin{Highlighting}[]
\NormalTok{fig, ax }\OperatorTok{=}\NormalTok{ plt.subplots()}

\NormalTok{t }\OperatorTok{=}\NormalTok{ np.linspace(}\DecValTok{0}\NormalTok{, }\DecValTok{10}\NormalTok{, }\DecValTok{100}\NormalTok{)}
\NormalTok{ax.plot(t, np.sin(t), label}\OperatorTok{=}\StringTok{\textquotesingle{}sin(t)\textquotesingle{}}\NormalTok{)}
\NormalTok{ax.plot(t, np.cos(t), label}\OperatorTok{=}\StringTok{\textquotesingle{}cos(t)\textquotesingle{}}\NormalTok{)}
\NormalTok{ax.plot(t, np.sin(t)}\OperatorTok{**}\DecValTok{2}\NormalTok{, label}\OperatorTok{=}\StringTok{\textquotesingle{}sin²(t)\textquotesingle{}}\NormalTok{)}

\NormalTok{ax.set\_xlabel(}\StringTok{\textquotesingle{}Time\textquotesingle{}}\NormalTok{)}
\NormalTok{ax.set\_ylabel(}\StringTok{\textquotesingle{}Value\textquotesingle{}}\NormalTok{)}
\NormalTok{ax.set\_title(}\StringTok{\textquotesingle{}Trigonometric Functions\textquotesingle{}}\NormalTok{)}
\NormalTok{ax.legend()}
\NormalTok{ax.grid(}\VariableTok{True}\NormalTok{, alpha}\OperatorTok{=}\FloatTok{0.3}\NormalTok{)}

\NormalTok{plt.show()}
\end{Highlighting}
\end{Shaded}

\textbf{Key elements:} labels, legend, grid for readability.

\subsection{Scatter Plots}\label{scatter-plots}

\textbf{For showing relationships:}

\begin{Shaded}
\begin{Highlighting}[]
\NormalTok{np.random.seed(}\DecValTok{42}\NormalTok{)}
\NormalTok{x }\OperatorTok{=}\NormalTok{ np.random.randn(}\DecValTok{100}\NormalTok{)}
\NormalTok{y }\OperatorTok{=} \DecValTok{2}\OperatorTok{*}\NormalTok{x }\OperatorTok{+}\NormalTok{ np.random.randn(}\DecValTok{100}\NormalTok{)}

\NormalTok{fig, ax }\OperatorTok{=}\NormalTok{ plt.subplots()}
\NormalTok{ax.scatter(x, y, alpha}\OperatorTok{=}\FloatTok{0.6}\NormalTok{, s}\OperatorTok{=}\DecValTok{50}\NormalTok{)}
\NormalTok{ax.set\_xlabel(}\StringTok{\textquotesingle{}x\textquotesingle{}}\NormalTok{)}
\NormalTok{ax.set\_ylabel(}\StringTok{\textquotesingle{}y\textquotesingle{}}\NormalTok{)}
\NormalTok{ax.set\_title(}\StringTok{\textquotesingle{}Scatter Plot Example\textquotesingle{}}\NormalTok{)}
\NormalTok{ax.grid(}\VariableTok{True}\NormalTok{, alpha}\OperatorTok{=}\FloatTok{0.3}\NormalTok{)}
\NormalTok{plt.show()}
\end{Highlighting}
\end{Shaded}

\textbf{Economic use:} Plotting data points, cross-country regressions,
identifying outliers.

\subsection{Customization: Colors and
Styles}\label{customization-colors-and-styles}

\textbf{Colors:}

\begin{Shaded}
\begin{Highlighting}[]
\NormalTok{ax.plot(x, y, color}\OperatorTok{=}\StringTok{\textquotesingle{}blue\textquotesingle{}}\NormalTok{)     }\CommentTok{\# Named colors}
\NormalTok{ax.plot(x, y, color}\OperatorTok{=}\StringTok{\textquotesingle{}\#FF5733\textquotesingle{}}\NormalTok{)  }\CommentTok{\# Hex codes}
\NormalTok{ax.plot(x, y, color}\OperatorTok{=}\NormalTok{(}\FloatTok{0.2}\NormalTok{, }\FloatTok{0.4}\NormalTok{, }\FloatTok{0.6}\NormalTok{))  }\CommentTok{\# RGB tuple}
\end{Highlighting}
\end{Shaded}

\textbf{Line styles:}

\begin{Shaded}
\begin{Highlighting}[]
\NormalTok{ax.plot(x, y, linestyle}\OperatorTok{=}\StringTok{\textquotesingle{}{-}\textquotesingle{}}\NormalTok{)   }\CommentTok{\# Solid (default)}
\NormalTok{ax.plot(x, y, linestyle}\OperatorTok{=}\StringTok{\textquotesingle{}{-}{-}\textquotesingle{}}\NormalTok{)  }\CommentTok{\# Dashed}
\NormalTok{ax.plot(x, y, linestyle}\OperatorTok{=}\StringTok{\textquotesingle{}:\textquotesingle{}}\NormalTok{)   }\CommentTok{\# Dotted}
\NormalTok{ax.plot(x, y, linestyle}\OperatorTok{=}\StringTok{\textquotesingle{}{-}.\textquotesingle{}}\NormalTok{)  }\CommentTok{\# Dash{-}dot}
\end{Highlighting}
\end{Shaded}

\textbf{Markers:}

\begin{Shaded}
\begin{Highlighting}[]
\NormalTok{ax.plot(x, y, marker}\OperatorTok{=}\StringTok{\textquotesingle{}o\textquotesingle{}}\NormalTok{)  }\CommentTok{\# Circles}
\NormalTok{ax.plot(x, y, marker}\OperatorTok{=}\StringTok{\textquotesingle{}s\textquotesingle{}}\NormalTok{)  }\CommentTok{\# Squares}
\end{Highlighting}
\end{Shaded}

\subsection{Subplots}\label{subplots}

\textbf{Multiple plots in one figure:}

\begin{Shaded}
\begin{Highlighting}[]
\NormalTok{fig, (ax1, ax2) }\OperatorTok{=}\NormalTok{ plt.subplots(}\DecValTok{1}\NormalTok{, }\DecValTok{2}\NormalTok{, figsize}\OperatorTok{=}\NormalTok{(}\DecValTok{10}\NormalTok{, }\DecValTok{4}\NormalTok{))}

\CommentTok{\# Left plot}
\NormalTok{ax1.plot(x, np.sin(x))}
\NormalTok{ax1.set\_title(}\StringTok{\textquotesingle{}Sine\textquotesingle{}}\NormalTok{)}

\CommentTok{\# Right plot}
\NormalTok{ax2.plot(x, np.cos(x))}
\NormalTok{ax2.set\_title(}\StringTok{\textquotesingle{}Cosine\textquotesingle{}}\NormalTok{)}

\NormalTok{plt.tight\_layout()  }\CommentTok{\# Prevents overlap}
\NormalTok{plt.show()}
\end{Highlighting}
\end{Shaded}

\textbf{Grid layouts:} \texttt{plt.subplots(2,\ 2)} for 2×2 grid.

\subsection{Economic Application: Time
Series}\label{economic-application-time-series}

\textbf{Plotting GDP growth:}

\begin{Shaded}
\begin{Highlighting}[]
\NormalTok{years }\OperatorTok{=}\NormalTok{ np.arange(}\DecValTok{1980}\NormalTok{, }\DecValTok{2024}\NormalTok{)}
\NormalTok{gdp\_growth }\OperatorTok{=} \DecValTok{3} \OperatorTok{+}\NormalTok{ np.random.randn(}\BuiltInTok{len}\NormalTok{(years)) }\OperatorTok{*} \DecValTok{2}

\NormalTok{fig, ax }\OperatorTok{=}\NormalTok{ plt.subplots(figsize}\OperatorTok{=}\NormalTok{(}\DecValTok{10}\NormalTok{, }\DecValTok{5}\NormalTok{))}
\NormalTok{ax.plot(years, gdp\_growth, linewidth}\OperatorTok{=}\DecValTok{2}\NormalTok{)}
\NormalTok{ax.axhline(y}\OperatorTok{=}\DecValTok{0}\NormalTok{, color}\OperatorTok{=}\StringTok{\textquotesingle{}r\textquotesingle{}}\NormalTok{, linestyle}\OperatorTok{=}\StringTok{\textquotesingle{}{-}{-}\textquotesingle{}}\NormalTok{, alpha}\OperatorTok{=}\FloatTok{0.5}\NormalTok{)}
\NormalTok{ax.fill\_between(years, }\DecValTok{0}\NormalTok{, gdp\_growth, }
\NormalTok{                 where}\OperatorTok{=}\NormalTok{(gdp\_growth }\OperatorTok{\textless{}} \DecValTok{0}\NormalTok{), }
\NormalTok{                 alpha}\OperatorTok{=}\FloatTok{0.3}\NormalTok{, color}\OperatorTok{=}\StringTok{\textquotesingle{}red\textquotesingle{}}\NormalTok{, label}\OperatorTok{=}\StringTok{\textquotesingle{}Recessions\textquotesingle{}}\NormalTok{)}

\NormalTok{ax.set\_xlabel(}\StringTok{\textquotesingle{}Year\textquotesingle{}}\NormalTok{)}
\NormalTok{ax.set\_ylabel(}\StringTok{\textquotesingle{}GDP Growth Rate (\%)\textquotesingle{}}\NormalTok{)}
\NormalTok{ax.set\_title(}\StringTok{\textquotesingle{}U.S. GDP Growth Rate\textquotesingle{}}\NormalTok{)}
\NormalTok{ax.legend()}
\NormalTok{ax.grid(}\VariableTok{True}\NormalTok{, alpha}\OperatorTok{=}\FloatTok{0.3}\NormalTok{)}
\NormalTok{plt.show()}
\end{Highlighting}
\end{Shaded}

\subsection{Economic Application:
Distributions}\label{economic-application-distributions}

\textbf{Histogram and kernel density:}

\begin{Shaded}
\begin{Highlighting}[]
\ImportTok{from}\NormalTok{ scipy.stats }\ImportTok{import}\NormalTok{ norm}

\NormalTok{income }\OperatorTok{=}\NormalTok{ np.random.lognormal(}\DecValTok{10}\NormalTok{, }\FloatTok{0.5}\NormalTok{, }\DecValTok{1000}\NormalTok{)}

\NormalTok{fig, (ax1, ax2) }\OperatorTok{=}\NormalTok{ plt.subplots(}\DecValTok{1}\NormalTok{, }\DecValTok{2}\NormalTok{, figsize}\OperatorTok{=}\NormalTok{(}\DecValTok{12}\NormalTok{, }\DecValTok{4}\NormalTok{))}

\CommentTok{\# Histogram}
\NormalTok{ax1.hist(income, bins}\OperatorTok{=}\DecValTok{30}\NormalTok{, edgecolor}\OperatorTok{=}\StringTok{\textquotesingle{}black\textquotesingle{}}\NormalTok{, alpha}\OperatorTok{=}\FloatTok{0.7}\NormalTok{)}
\NormalTok{ax1.set\_xlabel(}\StringTok{\textquotesingle{}Income\textquotesingle{}}\NormalTok{)}
\NormalTok{ax1.set\_ylabel(}\StringTok{\textquotesingle{}Frequency\textquotesingle{}}\NormalTok{)}
\NormalTok{ax1.set\_title(}\StringTok{\textquotesingle{}Income Distribution\textquotesingle{}}\NormalTok{)}

\CommentTok{\# Log scale to see tail}
\NormalTok{ax2.hist(income, bins}\OperatorTok{=}\DecValTok{30}\NormalTok{, edgecolor}\OperatorTok{=}\StringTok{\textquotesingle{}black\textquotesingle{}}\NormalTok{, alpha}\OperatorTok{=}\FloatTok{0.7}\NormalTok{)}
\NormalTok{ax2.set\_xlabel(}\StringTok{\textquotesingle{}Income\textquotesingle{}}\NormalTok{)}
\NormalTok{ax2.set\_yscale(}\StringTok{\textquotesingle{}log\textquotesingle{}}\NormalTok{)}
\NormalTok{ax2.set\_title(}\StringTok{\textquotesingle{}Income Distribution (Log Scale)\textquotesingle{}}\NormalTok{)}

\NormalTok{plt.tight\_layout()}
\NormalTok{plt.show()}
\end{Highlighting}
\end{Shaded}

\subsection{Economic Application: Supply and
Demand}\label{economic-application-supply-and-demand}

\begin{Shaded}
\begin{Highlighting}[]
\NormalTok{q }\OperatorTok{=}\NormalTok{ np.linspace(}\DecValTok{0}\NormalTok{, }\DecValTok{100}\NormalTok{, }\DecValTok{100}\NormalTok{)}
\NormalTok{demand }\OperatorTok{=} \DecValTok{100} \OperatorTok{{-}}\NormalTok{ q}
\NormalTok{supply }\OperatorTok{=}\NormalTok{ q}

\NormalTok{fig, ax }\OperatorTok{=}\NormalTok{ plt.subplots()}
\NormalTok{ax.plot(q, demand, label}\OperatorTok{=}\StringTok{\textquotesingle{}Demand\textquotesingle{}}\NormalTok{, linewidth}\OperatorTok{=}\DecValTok{2}\NormalTok{)}
\NormalTok{ax.plot(q, supply, label}\OperatorTok{=}\StringTok{\textquotesingle{}Supply\textquotesingle{}}\NormalTok{, linewidth}\OperatorTok{=}\DecValTok{2}\NormalTok{)}

\CommentTok{\# Equilibrium}
\NormalTok{q\_eq }\OperatorTok{=} \DecValTok{50}
\NormalTok{p\_eq }\OperatorTok{=} \DecValTok{50}
\NormalTok{ax.plot(q\_eq, p\_eq, }\StringTok{\textquotesingle{}ro\textquotesingle{}}\NormalTok{, markersize}\OperatorTok{=}\DecValTok{10}\NormalTok{, label}\OperatorTok{=}\StringTok{\textquotesingle{}Equilibrium\textquotesingle{}}\NormalTok{)}

\NormalTok{ax.set\_xlabel(}\StringTok{\textquotesingle{}Quantity\textquotesingle{}}\NormalTok{)}
\NormalTok{ax.set\_ylabel(}\StringTok{\textquotesingle{}Price\textquotesingle{}}\NormalTok{)}
\NormalTok{ax.set\_title(}\StringTok{\textquotesingle{}Supply and Demand\textquotesingle{}}\NormalTok{)}
\NormalTok{ax.legend()}
\NormalTok{ax.grid(}\VariableTok{True}\NormalTok{, alpha}\OperatorTok{=}\FloatTok{0.3}\NormalTok{)}
\NormalTok{plt.show()}
\end{Highlighting}
\end{Shaded}

\subsection{Publication-Quality
Figures}\label{publication-quality-figures}

\textbf{Key elements for professional plots:}

\begin{Shaded}
\begin{Highlighting}[]
\NormalTok{plt.rcParams[}\StringTok{\textquotesingle{}figure.figsize\textquotesingle{}}\NormalTok{] }\OperatorTok{=}\NormalTok{ (}\DecValTok{10}\NormalTok{, }\DecValTok{6}\NormalTok{)}
\NormalTok{plt.rcParams[}\StringTok{\textquotesingle{}font.size\textquotesingle{}}\NormalTok{] }\OperatorTok{=} \DecValTok{12}
\NormalTok{plt.rcParams[}\StringTok{\textquotesingle{}lines.linewidth\textquotesingle{}}\NormalTok{] }\OperatorTok{=} \DecValTok{2}

\NormalTok{fig, ax }\OperatorTok{=}\NormalTok{ plt.subplots()}
\NormalTok{ax.plot(x, y)}
\NormalTok{ax.set\_xlabel(}\StringTok{\textquotesingle{}x\textquotesingle{}}\NormalTok{, fontsize}\OperatorTok{=}\DecValTok{14}\NormalTok{)}
\NormalTok{ax.set\_ylabel(}\StringTok{\textquotesingle{}y\textquotesingle{}}\NormalTok{, fontsize}\OperatorTok{=}\DecValTok{14}\NormalTok{)}
\NormalTok{ax.set\_title(}\StringTok{\textquotesingle{}Professional Figure\textquotesingle{}}\NormalTok{, fontsize}\OperatorTok{=}\DecValTok{16}\NormalTok{)}

\CommentTok{\# Save with high DPI}
\NormalTok{plt.savefig(}\StringTok{\textquotesingle{}figure.pdf\textquotesingle{}}\NormalTok{, dpi}\OperatorTok{=}\DecValTok{300}\NormalTok{, bbox\_inches}\OperatorTok{=}\StringTok{\textquotesingle{}tight\textquotesingle{}}\NormalTok{)}
\NormalTok{plt.savefig(}\StringTok{\textquotesingle{}figure.png\textquotesingle{}}\NormalTok{, dpi}\OperatorTok{=}\DecValTok{300}\NormalTok{, bbox\_inches}\OperatorTok{=}\StringTok{\textquotesingle{}tight\textquotesingle{}}\NormalTok{)}
\end{Highlighting}
\end{Shaded}

\textbf{Format tips:} PDF for papers, PNG for presentations.

\subsection{Summary: Key Takeaways}\label{summary-key-takeaways}

\begin{enumerate}
\def\labelenumi{\arabic{enumi}.}
\tightlist
\item
  \textbf{Matplotlib} creates publication-quality plots for economic
  analysis
\item
  \textbf{Object-oriented approach} (\texttt{fig,\ ax}) provides better
  control
\item
  \textbf{Customization} enables professional-looking figures
\item
  \textbf{Subplots} combine multiple visualizations effectively
\item
  \textbf{Economic applications} include time series, distributions,
  equilibrium diagrams
\item
  \textbf{Save figures} in appropriate formats (PDF, PNG) with high
  resolution
\end{enumerate}

\textbf{Next:} SciPy for optimization and numerical methods in
economics.




\end{document}
